% !TeX root = ../main.tex
\section{Измерение амплитудно-частотной характеристики фильтра}
В ходе эксперимента выполнено измерение коэффициента передачи исследуемого фильтра с предварительной калибровкой сквозного тракта для устранения затухания в соединительных кабелях и неравномерности АЧХ измерительных приборов.

Коэффициент передачи фильтра на каждой частотной точке рассчитывался как разность показаний анализатора спектра:
\begin{equation}
    K_f(f) = P_{\text{изм}}(f) - P_{\text{кал}}(f).
\end{equation}

График экспериментальной АЧХ исследуемого фильтра нижних частот представлен на рисунке~\ref{fig:ach}.

\begin{figure}[H]
    \centering
    \includegraphics[width=0.75\linewidth]{АЧХ.png}
    \caption{Амплитудно-частотная характеристика (АЧХ) исследуемого фильтра нижних частот с контрольными уровнями затухания $-1\text{ дБ}$ и $-3\text{ дБ}$}
    \label{fig:ach}
\end{figure}

По полученной экспериментальной кривой определены ключевые частотные параметры фильтра. Полоса пропускания по уровню $-1\text{ дБ}$ составляет приблизительно $f_{-1\text{ дБ}} \approx 468\text{ МГц}$, где измеренное значение коэффициента передачи достигает $K_f \approx -2{,}96\text{ дБ}$.

Граничная частота среза по уровню половинной мощности зафиксирована на значении $f_{-3\text{ дБ}} \approx 527\text{ МГц}$, при этом уровень передачи составил $K_f \approx -4{,}96\text{ дБ}$.

В полосе заграждения частота гарантированного подавления по уровню $-40\text{ дБ}$ составляет $f_{-40\text{ дБ}} \approx 657\text{ МГц}$, где уровень передачи спадает практически до $-40\text{ дБ}$, обеспечивая глубокое подавление внеполосных составляющих.

На основе определенных граничных частот рассчитан коэффициент прямоугольности фильтра:
\begin{equation}
    K_{\text{п}} = \frac{f_{-40\text{ дБ}}}{f_{-3\text{ дБ}}} = \frac{657\text{ МГц}}{527\text{ МГц}} \approx 1{,}25.
\end{equation}
Полученное значение свидетельствует о высокой крутизне ската АЧХ в переходной зоне и хорошей избирательности исследуемого ФНЧ.
